\chapter*{Abstract}
Sickle cell disease remains a high-impact hemoglobinopathy, particularly in regions where specialized laboratories are scarce. Conventional diagnostics depend on costly instrumentation, chemical staining, and expert cytotechnologists, which hinders timely screening. This thesis introduces an automated system that classifies holograms captured with lensless digital holographic microscopy. The guiding hypothesis is that raw interference patterns contain enough discriminative cues to separate healthy erythrocytes from sickled cells when texture, morphology, and frequency descriptors are combined with classic interpretable machine-learning models and an anomaly detection layer.

The study is grounded on a curated dataset of 304 labeled holograms (151 healthy, 153 sickle cell) and an external repository of non-blood specimens. The methodology encompasses automatic dataset unification, extraction of 84 physics-inspired features, feature selection, and the construction of an ensemble that blends logistic regression, support vector machines, random forests, and gradient boosting classifiers. An anomaly detector based on the Mahalanobis distance complements the ensemble by flagging out-of-distribution samples and providing interpretable evidence for extreme cases.

The end-to-end pipeline is delivered as a reproducible command-line interface that supports quick, full, and deep analysis modes. Experimental campaigns reached a test accuracy of 95.1\% in quick mode and a ROC-AUC of 0.995 in deep mode, leveraging stratified cross-validation and Leave-One-Out schemes. The anomaly detector correctly identified 100\% of the non-blood holograms. These outcomes demonstrate that the synergy between computational optics and interpretable machine learning enables portable, transparent, and scalable diagnostic solutions.

\textbf{Keywords:} sickle cell disease, lensless digital holographic microscopy, interpretable machine learning, anomaly detection, pattern analysis.
